\documentclass[11pt,a4paper]{article}

\usepackage[utf8]{inputenc}
\usepackage[T1]{fontenc}
\usepackage{XCharter}
\usepackage[brazil]{babel}
\usepackage[margin=2cm]{geometry}
\usepackage{amsmath, amssymb}
\usepackage[xcharter,bigdelims,vvarbb]{newtxmath}
% newtxmath's operators font requests the fontspec family `XCharter(0)' in OT1;
% without these substitutions math digits and operator names silently fall back
% to Computer Modern (LaTeX Font Warning: Font shape `OT1/XCharter(0)/m/n' undefined).
\DeclareFontFamily{OT1}{XCharter(0)}{}
\DeclareFontShape{OT1}{XCharter(0)}{m}{n}{<->ssub * XCharter-TLF/m/n}{}
\DeclareFontShape{OT1}{XCharter(0)}{m}{it}{<->ssub * XCharter-TLF/m/it}{}
\DeclareFontShape{OT1}{XCharter(0)}{b}{n}{<->ssub * XCharter-TLF/b/n}{}
\DeclareFontShape{OT1}{XCharter(0)}{bx}{n}{<->ssub * XCharter-TLF/b/n}{}
\usepackage{enumerate}
\usepackage{array}
\usepackage{tikz}
\usetikzlibrary{positioning, matrix}

\pagestyle{empty}
\setlength{\parindent}{0pt}
\setlength{\parskip}{0.5em}

\renewcommand{\arraystretch}{1.4}

\begin{document}

%% =============================================================
%%  Header
%% =============================================================
{\Large\bfseries Aprendizado Profundo} \hfill \textbf{Prova, Unidade 4}\\
Prof. Lucas S. Kupssinskü \hfill Data: \rule{3cm}{0.4pt}\\[0.4em]
Nome: \rule{12cm}{0.4pt}

\rule{\linewidth}{0.8pt}

%% =============================================================
%%  Response grid
%% =============================================================
\textbf{Quadro de respostas (questões objetivas):}\\[0.3em]
\begin{center}
\begin{tabular}{|c|c|c|c|c|c|c|c|c|}
\hline
\textbf{Questão} & \textbf{a} & \textbf{b} & \textbf{c} & \textbf{d} & \textbf{e} & \textbf{f} & \textbf{g} & \textbf{h} \\\hline
1 & & & & & & & & \\\hline
2 & & & & & & & & \\\hline
3 & & & & & & & & \\\hline
4 & & & & & & & & \\\hline
5 & & & & & & & & \\\hline
\end{tabular}
\end{center}

\rule{\linewidth}{0.4pt}

%% =============================================================
%%  Objective questions  (5 x 1.2 = 6.0 pts)
%% =============================================================
\section*{Questões Objetivas (6.0 pts)}
\textit{Cada questão possui apenas uma alternativa correta. Marque a letra escolhida no quadro de respostas.}

\begin{enumerate}[1)]

%% ---------- Q1: OBJ-A numerical (KV-cache calculation across MHA / GQA) ----------
\item (1.2) Considere um modelo \textit{decoder-only} com $H = 8$ cabeças de \textit{query}, dimensão por cabeça $d_k = 64$, $L = 32$ camadas e \textit{batch size} $1$. Você usa \textit{KV-cache} durante a inferência. Quantos \textit{floats} por token compõem o \textit{KV-cache} para \textit{Multi-Head Attention} (MHA, $H$ cabeças de K e V) e para \textit{Grouped-Query Attention} (GQA com $G = 2$ grupos para K e V), e qual a razão entre as duas?
\begin{enumerate}[a)]
    \item MHA: $32\,768$; GQA: $16\,384$; GQA é $2\times$ menor que MHA.
    \item MHA: $32\,768$; GQA: $8\,192$; GQA é $4\times$ menor que MHA.
    \item MHA: $16\,384$; GQA: $4\,096$; GQA é $4\times$ menor que MHA.
    \item MHA: $32\,768$; GQA: $4\,096$; GQA é $8\times$ menor que MHA.
    \item MHA: $32\,768$; GQA: $2\,048$; GQA é $16\times$ menor que MHA.
    \item MHA: $16\,384$; GQA: $8\,192$; GQA é $2\times$ menor que MHA.
    \item MHA: $65\,536$; GQA: $32\,768$; GQA é $2\times$ menor que MHA.
    \item MHA e GQA armazenam o mesmo número de \textit{floats} por token, pois ambos cacheiam todas as cabeças de \textit{query}.
\end{enumerate}

\pagebreak

%% ---------- Q2: OBJ-B I/II/III analysis (tokenization) ----------
\item (1.2) Sobre as asserções abaixo a respeito de algoritmos de tokenização:

\textbf{I.} O \textbf{BPE} (\textit{Byte-Pair Encoding}) inicia com um vocabulário de caracteres (ou \textit{bytes}) e \emph{cresce} o vocabulário fundindo iterativamente o par de tokens mais \emph{frequente} no corpus de treinamento, até atingir o tamanho-alvo.

\textbf{II.} O \textbf{WordPiece} difere do BPE no critério de seleção: em vez do par mais frequente, mescla o par que \emph{maximiza a verossimilhança} dos dados de treinamento; assim como o BPE, parte de um vocabulário pequeno e cresce o vocabulário a cada iteração.

\textbf{III.} A \textbf{tokenização em nível de caractere} (cada letra/caractere vira um token) é, em geral, preferível aos métodos sub-palavra (BPE/WordPiece) em LLMs modernos, pois produz um vocabulário muito menor, elimina por completo o problema de \textit{out-of-vocabulary} e não acarreta custos adicionais relevantes de modelagem ou de tempo de treinamento.

\begin{enumerate}[a)]
    \item Apenas I está correta.
    \item Apenas II está correta.
    \item Apenas III está correta.
    \item Apenas I e II estão corretas.
    \item Apenas I e III estão corretas.
    \item Apenas II e III estão corretas.
    \item Todas estão corretas.
    \item Nenhuma está correta.
\end{enumerate}

%% ---------- Q3: OBJ-C scenario diagnostic (causal masking failure) ----------
\item (1.2) Você está pré-treinando um \textit{decoder-only} de 1B parâmetros com objetivo de \textit{next-token prediction} sobre um corpus de texto. Após implementar a sua versão de \textit{self-attention}, você observa que durante o treino o \textit{loss} cai rapidamente para próximo de zero, mas, em inferência, o modelo é incapaz de gerar texto coerente, produzindo apenas ruído. Qual a explicação mais consistente \textbf{e} a melhor mitigação?
\begin{enumerate}[a)]
    \item A taxa de aprendizado está alta demais, fazendo com que o \textit{loss} caia rápido demais; reduzi-la em duas ordens de grandeza resolve a divergência inferência-treino.
    \item O modelo está sofrendo \textit{overfitting}; reduzir o número de parâmetros e adicionar \textit{Dropout} resolve o problema.
    \item A máscara causal não está sendo aplicada a $QK^\top$, então o modelo enxerga tokens futuros e aprende a copiá-los; aplicar $-\infty$ às entradas $(i, j)$ com $j > i$ antes do \textit{softmax} resolve.
    \item O \textit{tokenizer} está mal treinado, fragmentando palavras em caracteres individuais; treiná-lo com mais corpus resolve o problema.
    \item Os \textit{positional encodings} estão sendo somados à entrada; removê-los faz com que o modelo deixe de depender de informação posicional e generalize melhor.
    \item O \textit{softmax} da atenção precisa ser substituído por \textit{sigmoid} ponto-a-ponto para problemas de geração autorregressiva.
    \item O modelo precisa de \textit{cross-attention}; \textit{decoder-only} sem \textit{cross-attention} é fundamentalmente incapaz de gerar texto coerente.
    \item O problema é que se trata de \textit{decoder-only}; trocar para \textit{encoder-only} (tipo BERT) resolve, pois BERT é o estado da arte em geração de texto.
\end{enumerate}

%% ---------- Q4: OBJ-D architectural-purpose recall (positional encoding) ----------
\item (1.2) Por que precisamos de \textit{positional encoding} em \textit{transformers}, enquanto RNNs não precisam dessa informação adicional?
\begin{enumerate}[a)]
    \item Porque sem \textit{positional encoding} o \textit{softmax} saturaria e o gradiente desapareceria; o \textit{positional encoding} ataca o problema de \textit{vanishing gradient}.
    \item Porque \textit{self-attention} tem complexidade $O(n^2)$ e o \textit{positional encoding} é o que reduz essa complexidade para aproximadamente $O(n)$.
    \item Porque a função \textit{softmax} usada na atenção é não-monótona, e o \textit{positional encoding} linearizá-a.
    \item Porque RNNs já incluem \textit{positional encoding} implícito via \textit{gates} de LSTM/GRU; \textit{transformers} omitem esse \textit{gating}.
    \item Porque \textit{self-attention} é permutação-equivariante: sem PE, ``o gato persegue o rato'' e ``o rato persegue o gato'' produzem as mesmas representações. RNNs codificam posição implicitamente pela ordem da recursão.
    \item Porque \textit{positional encoding} é necessário apenas durante a inferência; em treino, o gradiente fornece a informação posicional implicitamente.
    \item Porque \textit{transformers} usam apenas multiplicações e somas (sem recursão), e \textit{positional encoding} é a forma de introduzir não-linearidade na rede.
    \item Porque \textit{positional encoding} é uma técnica de regularização, equivalente em efeito ao \textit{Dropout} aplicado à camada de \textit{embeddings}.
\end{enumerate}

%% ---------- Q5: OBJ extra (causal mask pattern via TikZ visualization) ----------
\item (1.2) Os três padrões abaixo representam matrizes de atenção $A_{ij}$ de uma sequência de 4 tokens. Células \textbf{cinzas} indicam pares $(i, j)$ atendidos (peso possivelmente não-nulo após \textit{softmax}); células \textbf{brancas} indicam pares mascarados ($-\infty$ antes do \textit{softmax}, peso 0 depois).

\begin{center}
\begin{tikzpicture}[scale=0.45, every node/.style={font=\footnotesize}]
    %% Pattern A: lower-triangular (causal)
    \begin{scope}[xshift=0cm]
        \node at (2, 5) {\textbf{Padrão A}};
        \foreach \i in {0,1,2,3} {
            \foreach \j in {0,1,2,3} {
                \pgfmathparse{ifthenelse(\j<=\i, 1, 0)}
                \ifnum\pgfmathresult=1
                    \fill[gray!50] (\j, 3-\i) rectangle ++(1,1);
                \fi
                \draw (\j, 3-\i) rectangle ++(1,1);
            }
        }
        \node at (-0.6, 3.5) {\scriptsize $i{=}1$};
        \node at (-0.6, 0.5) {\scriptsize $i{=}4$};
        \node at (0.5, -0.6) {\scriptsize $j{=}1$};
        \node at (3.5, -0.6) {\scriptsize $j{=}4$};
    \end{scope}
    %% Pattern B: full matrix
    \begin{scope}[xshift=7cm]
        \node at (2, 5) {\textbf{Padrão B}};
        \foreach \i in {0,1,2,3} {
            \foreach \j in {0,1,2,3} {
                \fill[gray!50] (\j, 3-\i) rectangle ++(1,1);
                \draw (\j, 3-\i) rectangle ++(1,1);
            }
        }
        \node at (-0.6, 3.5) {\scriptsize $i{=}1$};
        \node at (-0.6, 0.5) {\scriptsize $i{=}4$};
        \node at (0.5, -0.6) {\scriptsize $j{=}1$};
        \node at (3.5, -0.6) {\scriptsize $j{=}4$};
    \end{scope}
    %% Pattern C: diagonal only
    \begin{scope}[xshift=14cm]
        \node at (2, 5) {\textbf{Padrão C}};
        \foreach \i in {0,1,2,3} {
            \foreach \j in {0,1,2,3} {
                \pgfmathparse{ifthenelse(\j==\i, 1, 0)}
                \ifnum\pgfmathresult=1
                    \fill[gray!50] (\j, 3-\i) rectangle ++(1,1);
                \fi
                \draw (\j, 3-\i) rectangle ++(1,1);
            }
        }
        \node at (-0.6, 3.5) {\scriptsize $i{=}1$};
        \node at (-0.6, 0.5) {\scriptsize $i{=}4$};
        \node at (0.5, -0.6) {\scriptsize $j{=}1$};
        \node at (3.5, -0.6) {\scriptsize $j{=}4$};
    \end{scope}
\end{tikzpicture}
\end{center}

Qual padrão é compatível com o pré-treinamento de um modelo \textit{decoder-only} autorregressivo (GPT, Llama)?
\begin{enumerate}[a)]
    \item Padrão A: atenção causal, necessária para \textit{next-token prediction} sem ver tokens futuros.
    \item Padrão A: ele permite que cada posição $i$ veja apenas a si mesma, garantindo independência entre os tokens durante o treino.
    \item Padrão B: \textit{decoder-only} requer atenção bidirecional para construir representações contextualizadas dos tokens passados.
    \item Padrão B: sem máscara o modelo aprende mais rapidamente; a propriedade autorregressiva é imposta apenas em inferência.
    \item Padrão C: cada token só atende a si mesmo, o que evita vazamento de informação futura.
    \item A e B são equivalentes do ponto de vista do treinamento autorregressivo.
    \item A e C bloqueiam tokens futuros, então ambos são corretos para \textit{decoder-only}.
    \item Nenhum: \textit{decoder-only} não usa atenção mascarada; apenas \textit{cross-attention} é mascarada.
\end{enumerate}

\end{enumerate}

\rule{\linewidth}{0.4pt}

%% =============================================================
%%  Dissertative questions  (4.0 pts)
%% =============================================================
\pagebreak
\section*{Questões Dissertativas (4.0 pts)}

\begin{enumerate}[1)]
\setcounter{enumi}{5}

%% ---------- Q6: DISS-A T/F with non-trivial modification ----------
\item (2.0) Marque \textbf{Verdadeiro} ou \textbf{Falso} nas asserções abaixo. Em caso \textbf{Falso}, proponha uma modificação \emph{não-trivial} que torna a asserção verdadeira (a simples negação \textbf{não} é uma modificação válida).

\begin{enumerate}[a)]
    \item ( )\,V \quad ( )\,F \quad Em \textit{Multi-Head Attention} com $H$ cabeças e dimensão $d_\text{model}$, cada cabeça opera com dimensão $d_k = d_\text{model}/H$ e a saída final é a concatenação das $H$ projeções, seguida por uma projeção linear $W^O$ que recompõe a dimensão $d_\text{model}$.\\[0.2em]
    \rule{\linewidth}{0.4pt}\\[0.2em]
    \rule{\linewidth}{0.4pt}

    \item ( )\,V \quad ( )\,F \quad O \textit{positional encoding} sinusoidal proposto por Vaswani et al.\ (2017) introduz, para cada uma das $d$ dimensões do \textit{embedding}, um vetor de parâmetros treináveis aprendidos por descida de gradiente durante o treinamento, o que permite à rede adaptar a representação posicional ao corpus específico.\\[0.2em]
    \rule{\linewidth}{0.4pt}\\[0.2em]
    \rule{\linewidth}{0.4pt}

    \item ( )\,V \quad ( )\,F \quad Modelos \textit{encoder-only} (e.g., BERT) usam atenção bidirecional e são tipicamente pré-treinados com objetivo de \textit{masked language modeling}; modelos \textit{decoder-only} (e.g., GPT) usam atenção mascarada (causal) e são pré-treinados com objetivo de \textit{next-token prediction}.\\[0.2em]
    \rule{\linewidth}{0.4pt}\\[0.2em]
    \rule{\linewidth}{0.4pt}

    \item ( )\,V \quad ( )\,F \quad O algoritmo BPE inicia com um vocabulário pequeno (caracteres ou \textit{bytes} individuais) e cresce iterativamente o vocabulário fundindo o par de tokens \emph{menos} frequente, até atingir o tamanho-alvo do vocabulário.\\[0.2em]
    \rule{\linewidth}{0.4pt}\\[0.2em]
    \rule{\linewidth}{0.4pt}

    \item ( )\,V \quad ( )\,F \quad \textit{Multi-Query Attention} (MQA) usa $H$ cabeças independentes para Q, K e V; portanto, reduzir o número de cabeças \emph{de Q} é a forma de diminuir o tamanho do \textit{KV-cache} durante a inferência.\\[0.2em]
    \rule{\linewidth}{0.4pt}\\[0.2em]
    \rule{\linewidth}{0.4pt}
\end{enumerate}

%% ---------- Q7: DISS-B numerical attention computation ----------
\item (2.0) Considere uma camada de \textit{self-attention} causal de \textbf{cabeça única} aplicada a uma sequência de 3 tokens $(x_1, x_2, x_3)$. As matrizes $Q$, $K$, $V$ já foram calculadas e têm dimensão $3 \times 1$ (sequência de comprimento 3, dimensão de chave $d_k = 1$):
\[
    Q = \begin{bmatrix} 1 \\ 1 \\ 1 \end{bmatrix},\quad
    K = \begin{bmatrix} 1 \\ 1 \\ 1 \end{bmatrix},\quad
    V = \begin{bmatrix} 2 \\ 4 \\ 6 \end{bmatrix}.
\]
A operação aplicada é $O = \mathrm{softmax}\!\left(\mathrm{mask}(QK^\top)/\sqrt{d_k}\right) V$, onde $\mathrm{mask}(\cdot)$ aplica a máscara causal.

\begin{enumerate}[a)]
    \item (0.5) Compute a matriz de pontuações $S = QK^\top$.
    \vspace{1.5cm}

    \item (0.5) Aplique a máscara causal a $S$ e indique a matriz $S_\text{mask}$ resultante.
    \vspace{1.5cm}

    \item (0.5) Aplique \textit{softmax} linha a linha em $S_\text{mask}$ e indique a matriz de pesos de atenção $A$.
    \vspace{2.5cm}

    \item (0.5) Compute a saída $O = A \cdot V$. Apresente os três valores $O_1$, $O_2$, $O_3$.
    \vspace{2cm}
\end{enumerate}

\end{enumerate}

\end{document}
